% Section 7: discussion. Evidence: VANET PAPER2_SCOPE_REVIEW.md §3 (surviving claim, circularity), §5 (R1/R2 assessment),
% §6 (required corrections -> reporting checklist); protocol §8.7, §8.10 finding 5, §8.11; plan §2 row 7 (limitations).
\section{Discussion}
\label{sec:discussion}

\subsection{A regulariser that preserves}
\label{sec:disc_preserves}

The one positive claim this paper makes for the method is that the decaying KL anchor prevents the collapse that the identically initialised, identically trained unanchored policy suffers. A reader may object that this is close to a tautology: a penalty towards a high-entropy distribution keeps entropy high. We think three things in the result are not given by the definition. The unanchored control collapses fast and far---per-state entropy over rollouts from the prior's 0.73 to 0.22--0.25 within 300 episodes, on both seeds---so the reward in this task does actively destroy the prior, and a warm start alone protects nothing. The training action-selection rule, which the first runs had confounded with the anchor, has no effect on either side. And the anchored policy is not a copy of the prior: its argmax composition changes, its lane oscillation disappears (340 to 10 transitions per episode), its agreement with human labels rises four points; the anchor leaves room for the policy to move, and the policy moves.

What the anchor does not do is make the policy better. On matched episodes under the execution rule it trains with, the frozen prior matches or beats the anchored policy on entropy, exits and speed, and proximity favours the anchored policy on one seed of two. The practical content of the result, for this task, is ``keep the cloned policy'': fine-tuning with the anchor returns approximately what went in, and fine-tuning without it returns something worse. That is the sense in which the anchor preserves and does not improve. Theorem~\ref{thm:kl} is consistent with it and adds nothing to it; with $\delta=0.999$ the anchor's weight falls only from 2.0 to 1.48 over the 300 episodes, so the regime in which the decay hands control to the reward was never reached, and this study says nothing about what the policy would become there.

Whether this is a negative result about the method or about the benchmark cannot be settled here. The measures on which improvement was sought---exits, speed, proximity---are the ones Section~\ref{sec:testability} shows to be solved by a constant policy, so a fine-tuned policy has nowhere to go but towards that constant, and under greedy execution it goes exactly there. A task in which the cloned prior is measurably deficient and the deficiency is not repaired by a constant would be a fair test of anchored fine-tuning; this one is not.

\subsection{What the prior and the execution rule decide}
\label{sec:disc_prior}

Two choices that the submitted paper did not treat as part of the method decide the behaviour more than the fine-tuning does. The first is the construction of the prior. Inverse-frequency class weights turn 0.26\,\% of lane-change labels into 22\,\% of the prior's argmax decisions on the validation split; the anchored policy's greedy action is then Lane\_Right in every state, which happens to be the optimal exit action in this network. Without the weights the prior's marginal is the human one, the anchored policy's greedy action is Maintain in every state, and under policy-attributable lane changes it exits in 7\,\% of cases. The apparent task competence of the submitted method was a by-product of a class-imbalance correction meeting a network with a one-line solution. The second is the execution rule. The same anchored checkpoint is a constant policy when executed greedily and a four-action stochastic policy when sampled, with exits of 100 against 83\,\% and speeds of 24.5 against 11.6\,m/s. A method that trains by sampling from a deliberately broad distribution and is then evaluated by argmax is being evaluated as a different policy from the one it trained, and its entropy and its safety figures are then properties of two different objects. Any method of this family should declare the execution rule as part of itself and report entropy and outcomes under that one rule.

\subsection{A reporting checklist}
\label{sec:disc_checklist}

The defects of Section~\ref{sec:defects} suggest a short list of things an evaluation of imitation-anchored RL in a traffic simulator should state, each of which would have exposed at least one of our errors had we been made to state it.
\begin{enumerate}
  \item \emph{One contract.} The observation features, their normalisation and the action order consumed by the prior, the fine-tuned actor, the regulariser's reference and every evaluator, recorded in the checkpoint and checked by the evaluators (D1--D3).
  \item \emph{Who moved the vehicle.} The simulator's lane-change and speed modes for controlled vehicles, and the count of manoeuvres the policy commanded beside the count the simulator made (D5).
  \item \emph{How many agents were counted.} Agents departed, tracked and scored per episode, and the definition of success per agent (D4).
  \item \emph{Reference policies that learn nothing.} Every constant-action policy that is plausibly competitive, and the simulator's default driver, through the same evaluator on the same episodes (D5, D6).
  \item \emph{Collisions as the simulator's.} Whether the simulator's driver model can produce a collision for a controlled vehicle at all; if not, no collision figure as a policy outcome (D6).
  \item \emph{Exposure.} Proximity or conflict counts per distance travelled as a pooled rate, not per episode (stranded and slow policies accumulate time without distance).
  \item \emph{A grouped split with logged access.} Recorded data split by vehicle or trajectory, with overlap checks, selection on validation only, and a log of every read of the test split (D7).
  \item \emph{The majority class.} The agreement of the constant majority policy, and of the prior itself, in every agreement table (D8).
  \item \emph{The entropy that the definition names.} Per-state entropy from the policy's own distribution, per actor before averaging, on recorded states and on visited states; a selected-action histogram only under that name (D9, D10).
  \item \emph{The execution rule.} Sampled or greedy, for training and for evaluation, stated with every figure and never mixed in a comparison.
  \item \emph{The frozen prior as a row.} The prior evaluated as a policy on the same episodes as the policy fine-tuned from it.
  \item \emph{Matched controls and a rule written first.} When a component's effect is claimed, a control that differs in that component alone, matched episodes, paired intervals over the independent unit, and thresholds fixed before the runs.
  \item \emph{Reward labelled.} Training reward as a training statistic, with its window and exploration setting, separate from anything the evaluator measures (D11).
\end{enumerate}
Items 2, 4, 5 and 11 are, in our experience, the least often reported, and they are the cheapest: two constant policies, the simulator's default and the frozen prior needed a few hundred evaluation episodes and no training, and they overturned most of the submitted safety table.

\subsection{Limitations}
\label{sec:disc_limits}

The isolation of the anchor rests on two training seeds per arm. The paired intervals quantify the evaluation uncertainty of those eight checkpoints over matched episodes and say nothing about variation across seeds; the collapse effect is large and of the same sign on both seeds and the frozen-prior comparison found no gain on either, but neither is a distributional statement. One scenario and one network were used, and Section~\ref{sec:testability} gives reasons to expect that network to flatter any policy that moves right. One prior construction, the class-balanced one, entered the isolation; the unweighted-prior ablation was trained under SUMO's default lane-change mode and the earlier coupling of sampling to the regulariser, so it shows what anchoring to a different prior produces but is not a matched arm of the design. The unanchored arms share the anchored arms' warm start, which is not a neutral initialisation for an unanchored learner (a randomly initialised MADDPG run from the earlier pilot did better than they did on every safety column), so the comparison isolates the anchor given the warm start and is not a comparison of CALM with MADDPG. The methods the submitted paper compared against---DAgger, GAIL, the soft actor--critic variant, the entropy-coefficient sweep and the $\lambda_0/\delta$ ablation---were not re-run under the corrected contract; QMIX-CALM was re-run on one seed for 100 episodes and is descriptive. The anchor's decay was not exercised within the training horizon. Reward is available only as a training statistic. One arm's per-episode entropies and lane transitions were lost to a logging overwrite and are reported without intervals (Table~\ref{tab:t6_frozen_prior}). All NGSIM figures are on the validation split, which shares its recording, and for most vehicles its time window, with the training split. And the defect catalogue is that of one pipeline, found by its own authors; we do not know how common each defect is in published work, only that none of them was visible in our reported numbers.
